\phantomsection\addcontentsline{toc}{section}{บทคัดย่อ}
\begin{center}{\bfseries\fontsize{20pt}{24pt}\selectfont บทคัดย่อ}
\end{center}
บทความนี้สรุปการออกแบบคลังข้อมูลการลานิสิตของคณะวิศวกรรมศาสตร์ กำแพงแสน โดยศึกษาซอร์สโค้ดส่วนจัดเตรียมข้อมูล บริการรายงาน และหน้ารายงานในโครงการ โครงสร้างหลักแยกข้อเท็จจริงระดับใบลาออกจากระดับใบลากับกลุ่มเรียน เชื่อมมิตินิสิต บุคลากร รายวิชา ภาคการศึกษา และหน่วยงาน เพื่อวิเคราะห์จำนวนใบลา วันลา และผลกระทบต่อการเรียนโดยลดการนับซ้ำ มุมมองรายงานครอบคลุมนิสิต อาจารย์ที่ปรึกษา อาจารย์ผู้สอน หัวหน้าภาควิชา ผู้บริหาร และเจ้าหน้าที่ โดยใช้ขอบเขตสิทธิ์ปัจจุบันร่วมกับบริบทข้อมูลในอดีต ผลการศึกษานี้เป็นการตรวจโครงสร้างและความสามารถจากโค้ด มิใช่ผลประเมินการใช้งานจริง เนื่องจากตัวจัดเตรียมข้อมูลหน้ารายงานยังใช้ข้อมูลจำลอง บทความจึงเสนอหลักการอ่านตัวชี้วัดและข้อจำกัดที่ต้องตรวจสอบก่อนนำข้อมูลไปใช้ตัดสินใจ

\noindent\textbf{คำสำคัญ:} คลังข้อมูล, การลานิสิต, แบบจำลองมิติ, รายงานตามบทบาท

\clearpage\phantomsection\addcontentsline{toc}{section}{Abstract}
\begin{center}{\bfseries\fontsize{20pt}{24pt}\selectfont Abstract}
\end{center}
This article examines the design of a student leave data warehouse for the Faculty of Engineering at Kamphaeng Saen. Source inspection covers data preparation, reporting services, and role-specific report pages. Separate request and request--offering facts support analysis of leave requests, calendar days, and teaching impacts while controlling duplicate aggregation. Shared dimensions describe students, personnel, subjects, academic terms, and organizations. Reporting perspectives cover students, advisors, lecturers, department heads, executives, and administrative staff. Current authorization scopes are distinguished from historical analytical context. Findings describe implemented structures and intended analytical capabilities rather than measured production outcomes: the report data provider currently supplies mock data. Metric definitions, data completeness, and integration limitations must therefore be verified before operational decision making.

\noindent\textbf{Keywords:} data warehouse, student leave, dimensional model, role-based reporting
